\documentclass[11pt,a4paper]{article}
\usepackage[margin=23mm]{geometry}
\usepackage{fontspec}\setmainfont{Latin Modern Roman}
\setmonofont{Latin Modern Mono}[Scale=0.85]
\usepackage{amsmath,amssymb,booktabs,longtable,xurl,hyperref}
\setlength{\parindent}{0pt}\setlength{\parskip}{6pt}
\newcommand{\code}[1]{\nolinkurl{#1}}
\title{S05 / REVIEW\_006\_001\\\large Przegląd wyników drugiego okna przed ustaleniem tezy}
\author{Przegląd własny z odrębnym replayem}\date{10 października 2026}
\begin{document}\maketitle
\textbf{Werdykt: conditional on a named input.}
Warunkowy rachunek MARGINAL-006 jest spójny w sprawdzonym zakresie.
Nie dostarcza efektywnego świadka dla all-h Q-JOINT.
Nie nadajemy projektowego REVIEWED, nie deklarujemy nowego kernela
ani pełnego niezależnego odbioru całego S05.

\section{Zakres i identyfikacja wyniku}
Sprawdzono zakończony wynik okna \emph{Sprawdź normalizator włókna QROM},
po commitach \code{f7588d0a}, \code{bb74c1a5}, \code{c957e798}.
Ostatni z nich jest bazą tego przeglądu. Odczyt stanu okna wykazał
zakończony turn i stan idle; nie wysłano tam wiadomości ani nowego zadania.

Zakres: dowód nowego marginalu, prawa capów, absolutnej ciągłości,
rachunku momentu i budżetu; korekta końców przedziałów w 005; wskazana
przeszkoda kosztowa i jej ograniczenia. Przeczytano źródła Sage oraz
istotne definicje użytych majorant. Nie jest to pełny audyt źródeł C,
Lean, wszystkich starszych pakietów ani całej literatury QROM.

Manifest MARGINAL-006: 18/18 pinów zgodnych.
Manifest skorygowanego NORMALIZER-005: 5/5 zgodnych.
Historia błędnych świadectw nie została nadpisana.

\section{Łańcuch zależności i sprawdzenie argumentu}
\begin{enumerate}
\item Pełnoboxowa propozycja $\nu$ i moneta $a$ mają mieć nowe, dwustronne
względne certyfikaty względem $U=w/G$ i $F/M$, gdzie $F=Z_1(c-hy)$.
Z dolnej granicy akceptacji i górnej granicy licznika wynika punktowa
dominacja zaakceptowanego marginalu. Przy dominacji kategorii $K_1$
otrzymujemy $Q_A\le dP_A$.
\item Niezależny produkt 16 par i ten sam kanał normy, cap16 i emisji
dają $Q\le CP$, $C=d^{16}$. Z normalizacji wynika $C\ge1$ oraz
$\sum Q^2/P\le C$. Nowy dowód nie dziedziczy certyfikatu h=0.
\item Każda z 16 par jest przygotowywana z własnym capem N.
Wyczerpanie dowolnego capu nadpisuje wynik końcowy porażką.
Prawo zaakceptowanej pary nie zależy od indeksu akceptacji, więc
$J=(1-\delta)Q+\delta\delta_\bot$ z
$\delta=1-(1-(1-\alpha)^N)^{16}$.
Ta definicja przygotowuje wszystkie 16 par. Leniwy algorytm kończący
pracę przy pierwszym sukcesie normy wymagałby odrębnego dowodu prawa.
\item Przy $p=P(\bot)>0$ obie składowe mieszanki są wspierane przez P.
Rozwinięcie kwadratu daje
\[
 \sum J^2/P=(1-\delta)^2\sum Q^2/P
  +2(1-\delta)\delta Q(\bot)/p+\delta^2/p.
\]
Nierówność trójkąta w $L^2(P)$ prowadzi do podanego boundu chi-kwadrat.
W przypadku $Q=P$ pozostaje dokładnie $1+\delta^2(1/p-1)$.
Przy $p=0$, $\delta>0$ nośnik rzeczywiście zostaje naruszony.
\end{enumerate}

Przypadki brzegowe $\delta=0$, $\delta=1$, $p=1$ zgadzają się z rachunkiem.
Nie zastąpiono małego TV małym momentem. Koszt $\delta^2/p$ jest zachowany,
a dodatniość p sama w sobie nie oznacza użytecznej dolnej granicy.

\small
\begin{longtable}{@{}p{.44\linewidth}p{.48\linewidth}@{}}
\toprule Obowiązek & Wynik przeglądu \\\midrule\endhead
Prawidłowe prawo po ważeniu & Passed dla zapisanych względnych przesłanek.\\
Prawo capu i wszystkie porażki & Passed dla konstrukcji z 16 przygotowaniami.\\
AC i moment mieszanki & Passed tekstowo pod $p>0$ i certyfikatem rdzenia.\\
Budżet $e<2^{-119}$ & Conditional: liczby policzone poprawnie pod nazwanymi
założeniami; brak ich wykonawczej instancji.\\
Koszt odrzucania całych wektorów & Passed w dokładnym zakresie h=1,c=0,
idealna propozycja U, stały envelope; nie ograniczenie wszystkich metod.\\
Pełnoboxowa propozycja i moneta & Open. Sampler 004 nie dostarcza
wymaganego dwustronnego certyfikatu przez samo posiadanie dominacji.\\
Wykonalność all-h & Open. Skończony cap $2^{336409}$ nie jest użyteczny.\\
Kernelizacja i source/quantum bridges & Not addressed przez ten przegląd.\\
\bottomrule\end{longtable}\normalsize

\section{Korekta normalizatora i kontrola rachunku}
Erratum poprawnie rozpoznaje materialny błąd: konwersja MPFR przez
\code{QQ(endpoint)} mogła przesunąć koniec przedziału do środka.
Dotyczyło to również składania obudów, nie tylko ich drukowania.
Nowe źródło używa \code{exact_rational()} i składa przedziały wyłącznie
z dokładnych końców wymiernych. W przeglądzie użyto wyłącznie
skorygowanego świadectwa jako wejścia numerycznego.

Własny świeży replay skorygowanego 005 w precyzji 512 odtworzył
\code{NORMALIZER_005_CORRECTED_CERTIFICATE.json} bajt w bajt.
Następnie replay MARGINAL-006 z tym świeżym wejściem odtworzył jego
512-bitowy certyfikat również bajt w bajt, w tym 102 dokładne prawa
testowe i deklarowane nierówności kosztu/budżetu.
To odrębne uruchomienie tej samej implementacji, nie niezależny algorytm
numeryczny ani dowód poprawności wszystkich bibliotek.

Dodatkowy własny checker \code{REVIEW_006_CAP_CHECK.sage} sprawdził
230 dokładnych praw na trzech atomach, w tym skrajne delty, $p=1$
i negatywną kontrolę nośnika. Sprawdza rozwinięcie momentu, konserwatywną
majorantę oraz przypadek dokładnego rdzenia. Te skończone kontrole
uzupełniają powyższy dowód, nie zastępują jego ogólnego argumentu.

Kosztowo istotny wynik $\alpha<2^{-767}$ dotyczy odrzucania \emph{całego}
wektora z U przy h=1,c=0. Nie wolno opisywać go jako dowodu niemożliwości
faktoryzacji blokowej lub jako ataku na bezpieczeństwo FT1536.
Analogicznie duży moment przed normą z 005 nie jest automatycznie
dużym końcowym momentem Q-JOINT.

\section{Co wolno teraz wykorzystać}
Można wykorzystać nowy warunkowy lemat mieszanki capu i dokładne wskazanie,
jak jego dodatkowy abort obciąża moment. Można odrzucić tę konkretną
strategię globalnego odrzucania jako praktycznego świadka zasobowego.
Nie można odhaczyć Q-JOINT, przyjąć małego e dla niewykonanej propozycji,
ani oznaczyć DyadicObstruction lub h=0 jako kernelowych.

Następny krok powinien mieć dwa równoległe kryteria: właściwe prawo
pełnej odpowiedzi oraz użyteczny koszt konstrukcji. Nowy kandydat musi
otrzymać własny certyfikat; kolejne zwiększenie precyzji nie usuwa
wykazanej małej akceptacji. Przed wyborem kandydata należy ustalić
główną tezę i rolę silniejszego all-h celu, jak w
\code{TARGET_CONTRACT_001.tex}.

Wyniki przeglądu i piny są zapisane osobno. Nie zmieniono bajtów
MARGINAL-006, NORMALIZER-005, kolekcji obstructions ani ich statusów.
\end{document}
